\documentclass[12pt]{article}

\usepackage{comment}
\usepackage{amsmath}
\usepackage{amssymb}  % for \nmid

\newcommand{\D}{{\mathbb D}}
\newcommand{\G}{{\mathbb G}}
\newcommand{\Z}{{\mathbb Z}}
\newcommand{\N}{{\mathbb N}}
\newcommand{\Q}{{\mathbb Q}}
\newcommand{\R}{{\mathbb R}}
\newcommand{\succc}{{\rm succ}}
\newcommand{\pred}{{\rm pred}}

\newcommand{\lc}{\left\lceil}
\newcommand{\rc}{\right\rceil}
\newcommand{\Ceil}[1]{\left\lceil {#1}\right\rceil}
\newcommand{\ceil}[1]{\left\lceil {#1}\right\rceil}
\newcommand{\floor}[1]{\left\lfloor{#1}\right\rfloor}

\begin{document}


\centerline{Homework 09, MORALLY Due April 21} 

\newif{\ifshowsoln}
\showsolntrue % comment out to NOT show solution inside \ifshowsoln and \fi blocks.

\begin{enumerate}
\item
(0 points)
What is your name.
\textbf{Gurkeerat Singh}

\vfill 

\centerline{\bf GO TO THE NEXT PAGE}

\newpage

\item 
(35 points) 

Read or re-read the slides on the Horse Numbers. 
The numbers $H(n)$ will come up in this problem.

For $n\ge 2$.
Let $I(n)$ be the number of ways that $n$ horses, $x_1,\ldots,x_n$, 

\noindent
can finish a race (so orderings with equalities allowed) that have $x_1<x_n$. 

\begin{enumerate}
\item
(0 points but you should do it or convince yourself that you could). 
What is $I(2)$, $I(3)$, $I(4)$. Do $I(4)$ in such a way that it can be
generalized to $I(n)$. 
\item
(30 points) 
Give a recurrence for $I(n)$. It may also involve $H(n)$. For example, it could be

$$I(n) = I(n-1)H(n-1) + 2I(n-2)H(n-2) + \cdots + (n-1)I(n-(n-1))H(n-(n-1)).$$

\boldmath{$\frac {H(n) - (H(n-1))}{2}$}

\end{enumerate}

\vfill 

\centerline{\bf GO TO THE NEXT PAGE}

\newpage

\item 
(35 points)

There are 30 people taking Game Theory.

There are 30 people taking Ramsey Theory.

There are 30 people taking Artificial Intelligence (AI).

There are 10 taking Game theory AND Ramsey Theory.

There are 10 taking Game theory AND AI.

There are 10 taking Ramsey Theory AND AI. 

There are 5 taking all three classes. 

\begin{enumerate}
\item 
How many students are there?\\ 
\textbf{We have to account for people that were multiple counted. So we have $30+30+30-10-10-10+5$ = 65 students}
\item
How many students are taking Game Theory AND Ramsey Theory but NOT AI?
\end{enumerate}
\textbf{We can split the 10 students taking Game Theory and Ramsey Theory into the 5 students taking all three classes and the 5 remaining students, which is the answer here. So, 5 Students.}

\vfill 

\centerline{\bf GO TO THE NEXT PAGE}

\newpage


\item 
(30 points-10 points each)

(This problem is based on material you will see Wednesday April 15.) 


Consider the equation: $w+x+y+z=100.$

\begin{enumerate}
\item 
How many solutions $(w,x,y,z)$ are there with $w,x,y,z\ge 0$?
\begin{itemize}
    \item{$103\choose 3$}
\end{itemize}

\item
How many solutions $(w,x,y,z)$ are there with $w,x,y,z\ge 5$?
\begin{itemize}
    \item{$83\choose 3$}
\end{itemize}
\item
How many solutions $(w,x,y,z)$ with $w\ge 1$ and $x\ge 10$? 
\begin{itemize}
    \item{$92\choose 3$}
\end{itemize}
\end{enumerate}

\end{enumerate}

\end{document}


